% Generated draft: requires mathematical and source review.
\documentclass[11pt]{article}
\usepackage[margin=25mm]{geometry}
\usepackage{fontspec}
\setmainfont{Latin Modern Roman}
\usepackage{amsmath,amssymb,mathtools,unicode-math}
\setmathfont{Latin Modern Math}
\usepackage{enumitem,needspace,xurl,hyperref}
\setlength{\parindent}{0pt}
\setlength{\parskip}{5pt}
\setlength{\emergencystretch}{3em}
\newcommand{\N}{\mathbb{N}}
\newcommand{\Z}{\mathbb{Z}}
\newcommand{\Q}{\mathbb{Q}}
\newcommand{\R}{\mathbb{R}}
\newcommand{\C}{\mathbb{C}}
\newcommand{\F}{\mathbb{F}}
\newcommand{\A}{\mathbb{A}}
\newcommand{\PP}{\mathbb{P}}
\newcommand{\E}{\mathbb{E}}
\newcommand{\eps}{\varepsilon}
\newcommand{\dd}{\,\mathrm{d}}
\begin{document}
\section*{UnsolvedMath ID: 30000721}
\textbf{Source problem code:} catalogue entry
\subsection{Definitions and mathematical statement}
For a simple graph $G$ and integer $k\geq1$, its power $G^k$ has vertex set $V(G)$, with distinct vertices adjacent exactly when their distance in $G$ is at most $k$. A graph is two-connected if it has at least three vertices, is connected, and remains connected after deletion of any one vertex. A ray is a one-way infinite simple path; two rays represent the same end if no finite vertex set separates their tails. Write $\Omega(H)$ for the ends of a graph $H$.

The source uses $|H|$ for a Freudenthal-type space obtained from the geometric graph by adjoining its ends. A Hamilton circle means an embedded copy of $S^1$ in this space containing every vertex and every end. In particular, it is a topological circle, not just a spanning double ray or an infinite closed walk.

The two separate questions are:
\begin{enumerate}
\item For every countable two-connected graph $G$, is there an embedding $c:S^1\to|G^2|$ whose image contains $V(G^2)\cup\Omega(G^2)$?
\item For every countable connected graph $G$, is there such an embedding into $|G^3|$, containing $V(G^3)\cup\Omega(G^3)$?
\end{enumerate}
The second question does not assume two-connectivity. Neither question is restricted to locally finite graphs: their extension beyond that setting is the supplied research target.

For connected locally finite graphs, the usual Freudenthal compactification supplies the topology and is compact Hausdorff. For arbitrary countable graphs of infinite vertex degree, the graph need not be locally compact and the same classical compactification does not automatically specify the intended space; an end topology and treatment of dominated ends must be supplied. The record does not provide these conventions. Its title indicates infinite graphs, but ``countable'' can include finite graphs: in the cube question, one- and two-vertex connected graphs have no topological Hamilton circle. The intended countably infinite or small-finite-exception scope must therefore be confirmed, not silently inferred as a theorem for every finite size.
\subsection{Short English statement}
Ask whether countable two-connected graph squares and countable connected graph cubes have circles spanning every vertex and end. The intended extension to non-locally-finite graphs requires a precise end-space topology, and the countable versus countably-infinite scope needs clarification.
\subsection{Sources}
\begin{itemize}
\item Supplied report: \emph{Graph Theory} (2008), Oberwolfach Reports, DOI 10.4171/owr/2007/16, pp. 887--944; proposed year 2007. \url{https://doi.org/10.4171/owr/2007/16}.
\item Supplied reference labelled ``infinite Hamilton circles'': \url{https://www.math.uni-hamburg.de/home/georgakopoulos/infinitefleischner.pdf}.
\item Local UnsolvedMath registry, permanent ID 30000721, catalogue entry; created 25 August 2026, updated 16 September 2026. Its corrected\_verified note says both questions and graph powers are explicit. The supplied 21 August 2026 literature triage labels the countable non-locally-finite conjectures open and distinguishes them from corresponding locally finite theorems. No PDF access or independent updated assessment is asserted here.
\end{itemize}
\end{document}
